\documentclass[12 pt, letterpaper]{hmcpset}
\usepackage{hw-preamble}

\name{Jayden Thadani}
\class{MATH 135 --- Complex Analysis}
\assignment{Homework 5}
\duedate{10th October 2024}

\begin{document}

\begin{problem}[14.]
	Prove that all entire functions that are also injective take the form $f(z) = az + b$ with $a, b \in \mathbb{C}$ and $a \neq 0$.
\end{problem}

\begin{solution}
	Suppose $f$ is entire and injective. Then by the chain rule, $g : z \mapsto f(1/z)$ is holomorphic on $\mathbb{C} \setminus \{ 0 \}$, and since $1/z$ is injective, so is $g$. Using the power series expansion for $f$, we can write a Laurent series expansion for $g$ on all of $\mathbb{C} \setminus \{ 0 \}$ --- if $f(z) = \sum_{n = 0}^{\infty} a_{n} z^{n}$, then
	\[
		g(z) = \sum_{n = 0}^{\infty} \frac{a_{n}}{z^{n}} = \sum_{n = -\infty}^{0} a_{-n} z^{n}.
	\]

	We claim that $g$ can't have infinitely many non-zero coefficients --- $g$ cannot have an essential singularity at $0$. If $g$ did, then by Casorati-Weierstrass, the images $g^{\text{img}}(\mathbb{D})$ and $g^{\text{img}}(\mathbb{C})$ would both be dense in $\mathbb{C}$, and then we would have some $g(a) = g(b)$ for $a \in \mathbb{D}$ and $b \in \mathbb{\mathbb{C}} \setminus \mathbb{D}$.

	If $g$ has only finitely many coefficients (suppose all $a_{-m}$ with $m > n$ are zero) terms in its Laurent series expansion, then $f$ must be a polynomial $a_{0} + a_{1} z + \dots + a_{n} z^{n}$. Suppose by way of contradiction that $n > 1$. Then, by the fundamental theorem of algebra, $f$ has $n$ roots. If any of them are distinct, then $f$ cannot be injective and we have a contradiction. If $f$ has $n$ repeated roots, we have $f(z) = (z - z_{0})^{n}$. Since $(z - z_{0})^{n} = 1$ has $n$ distinct solutions, $f$ is not injective here either, and again, we have a contradiction. Thus $n \le 1$. Of course, if $n < 1$ $f$ is just constant and not injective, so we must have $n = 1$. That is, $f = az + b$ with $a \neq 0$.
\end{solution}

\pagebreak

\begin{problem}[15.]
	Use the Cauchy inequalities or the maximum modulus principle to solve the following problems:
	\begin{enumerate}
		\item Prove that if $f$ is an entire function that satisfies
			\[
				\sup_{\lvert z \rvert = R} \lvert f(z) \rvert \le A R^{k} + B
			\]
			for all $R > 0$, and for some integer $k \ge 0$, and some constants $A, B > 0$, then $f$ is a polynomial of degree $\le k$.
		\setcounter{enumi}{3}
		\item Show that if the real part of an entire function $f$ is bounded, then $f$ is constant.
	\end{enumerate}
\end{problem}

\pagebreak

\begin{problem}[16.]
	Suppose $f$ and $g$ are holomorphic in a region containing the disc $\lvert z \rvert \le 1$. Suppose that $f$ has a simple zero at $z = 0$ and vanishes nowhere else in $\lvert z \rvert \le 1$.

	Let
	\[
		f_{\varepsilon}(z) = f(z) + \varepsilon g(z).
	\]
	Show that if $\varepsilon$ is sufficiently small, then
	\begin{enumerate}
		\item $f_{\varepsilon}(z)$ has a unique zero in $\lvert z \rvert \le 1$, and
		\item if $z_{\varepsilon}$ is this zero, the mapping $\varepsilon \mapsto z_{\varepsilon}$ is continuous.
	\end{enumerate}
\end{problem}

\pagebreak

\begin{problem}[18.]
	Give another proof of the Cauchy integral formula
	\[
		f(z) = \frac{1}{2 \pi i} \int_{C} \frac{f(\zeta)}{\zeta - z} \, d\zeta
	\]
	using homotopy of curves.
\end{problem}

\begin{solution}
	Suppose $f$ is holomorphic on the circle $C$ and the disc $D$ that it bounds. Then the quotient $\frac{f(\zeta)}{\zeta - z}$ is holomorphic on $C$ and $D$ except $z$. Notice that $C$ and the circle of radius $\varepsilon$ centred at $z$, $C_{\varepsilon}(z)$, are homotopic. Thus, 
	\[
		\int_{C} \frac{f(\zeta)}{\zeta - z} \, d\zeta = \int_{C_{\varepsilon}(z)} \frac{f(\zeta)}{\zeta - z} \, d\zeta.
	\]

	Now we can just imitate the standard proof of the Cauchy integral formula. We rewrite
	\[
		\int_{C_{\varepsilon}(z)} \frac{f(\zeta)}{\zeta - z} \, d\zeta = \int_{C_{\varepsilon}(z)} \frac{f(\zeta) - f(z)}{\zeta - z} \, d\zeta + f(z) \int_{C_{\varepsilon}(z)} \frac{1}{\zeta - z} \, d\zeta 
	\]
	Since $f$ is holomorphic, the first integrand, the difference quotient, is bounded, say by $M$, and thus
	\[
		\left \lvert \int_{C_{\varepsilon}(z)} \frac{f(\zeta) - f(z)}{\zeta - z} \, d \zeta  \right \rvert \le M 2 \pi \varepsilon
	\]
	Thus, the integral goes to zero as $\varepsilon \to 0$. The second integral, we've already verified is $2 \pi i$, and thus,
	\[
		\int_{C} \frac{f(\zeta)}{\zeta - z} \, d\zeta = \int_{C_{\varepsilon}(z)} \frac{f(\zeta)}{\zeta - z} \, d\zeta = 2 \pi i f(z)
	\]
	which is just the Cauchy integral formula, as desired.
\end{solution}

\pagebreak

\begin{problem}[19.]
	Prove the maximum principle for harmonic functions:
	\begin{enumerate}
		\item If $u$ is a non-constant real-valued harmonic function in a region $\Omega$, then $u$ cannot attain a maximum (or a minimum) in $\Omega$.
		\item Suppose that $\Omega$ is a region with compact closure $\overline{\Omega}$. If $u$ is harmonic in $\Omega$ and continuous in $\overline{\Omega}$, then
			\[
				\sup_{z \in \Omega} \lvert u(z) \rvert \le \sup_{z \in \overline{\Omega} \setminus \Omega} \lvert u(z) \rvert
			\]
	\end{enumerate}
\end{problem}

\end{document}
